
When an LLM generates code that fails tests, developers reach for two distinct feedback signals: static analysis errors from tools like pylint and mypy, and test failure messages from execution. Whether these signals detect the same bugs or fundamentally different ones has practical implications: if the signals are redundant, combining them wastes compute; if they are orthogonal, omitting either leaves bugs unfixed.

Recent work demonstrates that both feedback types independently improve LLM code generation quality. Static analysis feedback can reduce security issues from 40\% to 13\% over iterative refinement cycles. Execution-based feedback powers approaches like Self-Refine and CodeRL. However, no prior work quantifies whether static and execution feedback detect overlapping or distinct error classes.

This gap matters because the composition of feedback signals determines optimal refinement strategies. If static analysis catches a subset of what execution catches (or vice versa), then one signal suffices. If they catch disjoint error classes, then combining them should yield additive benefits.

We hypothesize that static analysis and execution feedback operate on fundamentally different code properties---structure versus behavior---and therefore detect categorically distinct error classes. Static tools analyze code text to find type mismatches, undefined variables, and syntax issues. Test suites execute code with inputs to find wrong outputs, runtime exceptions, and edge case failures.

We analyze 542 problems from HumanEval+ and MBPP+ to measure the Jaccard similarity between static-detected and execution-detected error sets. Our primary findings are:

\begin{enumerate}
\item \textbf{Complete orthogonality:} Jaccard similarity = 0.0 between static and execution error sets, indicating no overlap between detected error classes.
\end{enumerate}

\begin{enumerate}
\item \textbf{High static coverage:} Static analysis (pylint/mypy) detects errors in 98.6\% of code that fails execution tests, demonstrating its reliability as a feedback signal.
\end{enumerate}

\begin{enumerate}
\item \textbf{Behavioral detection not validated:} Only 0.3\% of static-clean code failed execution tests, far below the 40\% threshold. This result reflects the use of canonical solutions rather than actual LLM-generated code.
\end{enumerate}

